% DOCUMENT CLASS SETUP
\documentclass[12pt, a4paper]{article}

% DOCUMENT PACKAGES
\usepackage[a4paper, margin=1in]{geometry}
\usepackage{titling}
\usepackage{enumitem}
\usepackage{amsmath}
\usepackage{amssymb}
\usepackage{mathtools}
\usepackage{graphicx}
\usepackage{xcolor}
\usepackage{listings}
\usepackage{hyperref}
\usepackage{csquotes}
\usepackage{tocloft}

% FONT AND ENCODING SETUP
\usepackage{fontspec}
\usepackage{polyglossia}
\usepackage[Thai]{ucharclasses}
\setdefaultlanguage{english}
\setotherlanguage{thai}
\setTransitionsFor{Thai}{\thaifont}{\normalfont}
\newfontfamily{\thaifont}{THSarabun.ttf}[
    Path=./,
    BoldFont=THSarabunBold.ttf,
    ItalicFont=THSarabunItalic.ttf,
    BoldItalicFont=THSarabunBoldItalic.ttf,
    Scale=1.23
]
\XeTeXlinebreaklocale "th"
\XeTeXlinebreakskip = 0pt plus 1pt

% PROFESSIONAL CODE BLOCK
\usepackage{tcolorbox}
\tcbuselibrary{minted, skins, breakable}

% DEFINE GITHUB DARK THEME COLORS %
\definecolor{ghback}{HTML}{0d1117}    % Dark background
\definecolor{ghtext}{HTML}{c9d1d9}    % Light gray text
\definecolor{ghborder}{HTML}{30363d}  % GitHub Border color

% CODE BLOCK CONFIGURATION (MINTED) %
\renewcommand{\theFancyVerbLine}{\textcolor{white!60!gray}{\scriptsize\arabic{FancyVerbLine}}}
\newtcblisting{codeblock}[2][]{
    enhanced,
    breakable,
    colback=ghback,
    colframe=ghborder,
    coltitle=ghtext,
    colupper=ghtext,
    fonttitle=\bfseries\small,
    arc=3mm,
    boxrule=0.5mm,
    % PADDING SETTINGS
    left=25pt,
    right=10pt,
    top=5pt,
    bottom=5pt,
    % MINTED SETTINGS
    listing engine=minted,
    listing only,
    minted language=python,
    minted options={
        linenos,
        numbersep=5pt,
        fontsize=\footnotesize,
        breaklines=true,
        breakanywhere=true,
        autogobble,
        style=monokai,
    },
    title={#2},
    #1
}

% DOCUMENT TITLE SETUP %
\title {
	\textbf{2110446 DATA SCIENCE AND DATA ENGINEERING} \\
    \vspace{1em}
	\Large{\textbf{Homework 3: Decision Tree}} \\
    \vspace{0.5em}
    \large{Computer Engineering, Chulalongkorn University}
}
\author{Worralop Srichainont}
\date{\today}
\setlength{\droptitle}{15em}

% DISABLE PAGE NUMBER IN TABLE OF CONTENTS
\tocloftpagestyle{empty}

% DOCUMENT CONTENTS
\begin{document}

    \maketitle
    \thispagestyle{empty}
    \clearpage

    \tableofcontents
    \thispagestyle{empty}
    \clearpage

    \pagenumbering{arabic}

    \section{Section I}
        \subsection{Subsection I}
            \textbf{Kyphosis} is a condition where your spine curves outward more than it should. This causes your upper back around the thoracic region (the part of your spine between your neck and ribs) to bend forward.
        
        \subsection{Subsection II}
            The dataset has 81 rows and 4 columns which have the following details:
            \begin{enumerate}
                \item \textbf{\texttt{Kyphosis}} is the target column.
                \begin{enumerate}
                    \item \textbf{\texttt{absent}} means the spine has not deformed.
                    \item \textbf{\texttt{present}} means the spine has deformed.
                \end{enumerate}

                \item \textbf{\texttt{Age}} is the age of the patient in months unit.
                \item \textbf{\texttt{Number}} is the total number of vertebrae (spine bones) that were involved in the surgical procedure.
                \item \textbf{\texttt{Start}} is the index of the topmost vertebra (spine bone) that was operated on.
            \end{enumerate}

    \pagebreak

    \section{Section II}
        \subsection{Subsection I}
            To access the colab notebook, please follow this link: \texttt{\url{https://colab.research.google.com/drive/19gwAoeK6NIPp-c9uK-8UBJO1VauI06PL}}

        \subsection{Subsection II}
            Before we start, please import these libraries.

            \begin{codeblock}{Dependencies}
                import pandas as pd
                import matplotlib.pyplot as plt
                import seaborn as sns
                from sklearn.model_selection import train_test_split
                from sklearn.tree import DecisionTreeClassifier, plot_tree
                from sklearn.ensemble import RandomForestClassifier
                from sklearn.metrics import ConfusionMatrixDisplay, classification_report
            \end{codeblock}

    \pagebreak

    \section{Section III}
        \begin{table}[h]
            \centering
            \begin{tabular}{|c|c|c|c|c|}
                \hline
                \rule{0pt}{2.5ex}
                \textbf{Class} & \textbf{Precision} & \textbf{Recall} & \textbf{F1 Score} & \textbf{Support} \\
                \hline \hline
                \textbf{\texttt{absent}} & 0.9474 & 0.9000 & 0.9231 & 20 \\
                \textbf{\texttt{present}} & 0.6667 & 0.8000 & 0.7273 & 5 \\
                \hline
            \end{tabular}
            \caption{Class Metrics}
        \end{table}

        \begin{table}[h]
            \centering
            \begin{tabular}{|c|c|c|}
                \hline
                \rule{0pt}{2.5ex}
                \textbf{Accuracy} & \textbf{Macro F1 Score} & \textbf{Weighted F1 Score} \\
                \hline \hline
                0.8800 & 0.8252 & 0.8839 \\
                \hline
            \end{tabular}
            \caption{Overall Metrics}
        \end{table}

        \noindent \textbf{Result}: This decision tree model got \textbf{0.8252 F1 Macro Score}.

    \pagebreak

    \section{การประยุกต์ใช้การอินทิเกรต (Integration Application)}
        หากกำหนดอัตราการไหลของอากาศเข้าสู่ปอด ซึ่งเขียนด้วยฟังก์ชัน $f(t) = \frac{1}{2}\sin\left( \frac{2 \pi t}{5} \right)$ ในหน่วยลิตรต่อวินาที

        \vspace{0.5em}

        \noindent เราจะสามารถเขียนปริมาณของอากาศภายในปอด ในช่วงวินาทีที่ 0 ถึงวินาทีที่ t ได้ดังนี้ $V(t) = \int_{0}^{t} f(t) \,dt$

        \vspace{0.5em}

        \noindent กำหนดให้ $u = \frac{2\pi}{5}t$ จะได้ว่า $\,du = \frac{2\pi}{5}\,dt$
        \begin{align*}
            \int_{0}^{t} f(t) \,dt &= \int_{0}^{t} \frac{1}{2} \sin \left( \frac{2\pi}{5}t \right) \,dt \\[1em]
            \int_{0}^{t} f(t) \,dt &= \frac{5}{4\pi} \int_{0}^{t} \sin \left( \frac{2\pi}{5}t \right) \cdot \frac{2\pi}{5} \,dt \\[1em]
            \int_{0}^{t} f(t) \,dt &= \frac{5}{4\pi} \int_{0}^{t} \sin(u) \,du \\[1em]
            \int_{0}^{t} f(t) \,dt &= \frac{5}{4\pi} \int_{0}^{\frac{2\pi}{5}t} \sin(u) \,du \\[1em]
            \int_{0}^{t} f(t) \,dt &= \frac{5}{4\pi} \left[ -\cos(u) \right]_{0}^{\frac{2\pi}{5}t} \\[1em]
            \int_{0}^{t} f(t) \,dt &= \frac{5}{4\pi} \left[ -\cos\left( \frac{2\pi}{5}t \right) - (-\cos(0)) \right] \\[1em]
            \int_{0}^{t} f(t) \,dt &= \frac{5}{4\pi} \left[ 1 -\cos\left( \frac{2\pi}{5}t \right) \right]
        \end{align*}
        \noindent \textbf{ดังนั้น} จึงสามารถเขียนฟังก์ชันปริมาณของอากาศภายในปอด ณ วินาทีที่ t ได้เป็น $V(t) = \frac{5}{4\pi} \left[ 1 -\cos\left( \frac{2\pi}{5}t \right) \right]$

        \vspace{0.5em}

        \noindent ตัวอักษรปกติ \textbf{ตัวอักษรหนา} \textit{ตัวอักษรเอียง} \textbf{\textit{ตัวอักษรหนาและเอียง}} \underline{ขีดเส้นใต้}

        \noindent \enquote{ ข้อความแบบมีอัญประกาศ} \enquote{Quote text}
\end{document}
